\chapter*{Cómo usar esta guía}
\addcontentsline{toc}{chapter}{Cómo usar esta guía}

Esta guía es un repaso profundo, ordenado de lo más básico a lo más avanzado. Cada capítulo se apoya en los anteriores, así que
conviene leerla en orden la primera vez. Después puede usarse como consulta.

\begin{itemize}
  \item La \textbf{Parte I} construye las herramientas matemáticas: números complejos, fasores, impedancia, potencia,
        resonancia y decibelios.
  \item La \textbf{Parte II} desarrolla la teoría de radiofrecuencia exactamente como la implementa el motor: líneas,
        reflexión, carta de Smith, adaptación, parámetros S, estabilidad, modo mixto, dominio del tiempo, calibración, antenas y microstrip.
  \item La \textbf{Parte III} explica cómo está construida la aplicación, cómo se usa cada módulo y cómo se mantiene y publica.
\end{itemize}

Cada ecuación importante va seguida de un bloque \emph{«donde:»} que explica qué significa cada símbolo y en qué unidades se mide.
Los cuadros tienen un significado fijo:

\begin{definicion}{qué es}Definiciones formales de los conceptos.\end{definicion}
\begin{intuicion}[: qué significa]Explicación del concepto sin fórmulas: qué ocurre físicamente y por qué importa.\end{intuicion}
\begin{ejemplo}{con números}Ejercicios resueltos paso a paso. Los valores los ha calculado el motor de la aplicación.\end{ejemplo}
\begin{atencion}{errores frecuentes}Convenciones, trampas y confusiones habituales.\end{atencion}
\begin{enlaapp}Dónde aparece el concepto en la aplicación y cómo usarlo.\end{enlaapp}

% =====================================================================
\chapter*{Nomenclatura}
\addcontentsline{toc}{chapter}{Nomenclatura}

\section*{Matemática general}
\begin{longtable}{@{}>{\centering\arraybackslash}p{2.6cm}p{9.6cm}p{2.4cm}@{}}
\toprule \textbf{Símbolo} & \textbf{Significado} & \textbf{Unidad}\\ \midrule\endhead
$\jj$ & Unidad imaginaria, $\jj^2=-1$ (en ingeniería no se usa $i$ para no confundirla con la corriente) & --- \\
$z = a + \jj b$ & Número complejo genérico: $a$ parte real, $b$ parte imaginaria & --- \\
$\abs{z}$, $\varphi$ & Módulo (longitud) y argumento (ángulo) de un número complejo & ---, rad o $^\circ$ \\
$\conj{z}$ & Conjugado: cambia el signo de la parte imaginaria & --- \\
$\Re\{\cdot\}$, $\Im\{\cdot\}$ & Parte real y parte imaginaria & --- \\
$\ee^{\jj\varphi}$ & Exponencial compleja (fórmula de Euler) & --- \\
$\angle$ & «con ángulo»: $10\angle 30^\circ$ es un número de módulo 10 y ángulo $30^\circ$ & --- \\
$\ln$, $\log_{10}$ & Logaritmo natural (base $\ee$) y decimal & --- \\
$\mathcal{F}$, $\mathcal{F}^{-1}$ & Transformada de Fourier y su inversa & --- \\
\bottomrule
\end{longtable}

\section*{Señales y circuitos}
\begin{longtable}{@{}>{\centering\arraybackslash}p{2.6cm}p{9.6cm}p{2.4cm}@{}}
\toprule \textbf{Símbolo} & \textbf{Significado} & \textbf{Unidad}\\ \midrule\endhead
$t$ & Tiempo & \si{\second} \\
$f$ & Frecuencia: número de oscilaciones por segundo & \si{\hertz} \\
$\omega = 2\pi f$ & Frecuencia angular: radianes recorridos por segundo & \si{\radian\per\second} \\
$T = 1/f$ & Periodo: duración de una oscilación & \si{\second} \\
$v(t)$, $i(t)$ & Tensión y corriente instantáneas (valores reales en el tiempo) & \si{\volt}, \si{\ampere} \\
$V$, $I$ & Fasores de tensión y corriente (números complejos) & \si{\volt}, \si{\ampere} \\
$R$ & Resistencia: se opone a la corriente disipando energía & \si{\ohm} \\
$L$ & Inductancia: almacena energía en el campo magnético & \si{\henry} \\
$C$ & Capacidad: almacena energía en el campo eléctrico & \si{\farad} \\
$Z = R + \jj X$ & Impedancia: oposición compleja al paso de corriente alterna & \si{\ohm} \\
$X$ & Reactancia: parte imaginaria de $Z$ (positiva inductiva, negativa capacitiva) & \si{\ohm} \\
$X_L$, $X_C$ & Reactancia de una bobina y de un condensador & \si{\ohm} \\
$Y = G + \jj B$ & Admitancia, inversa de la impedancia & \si{\siemens} \\
$G$, $B$ & Conductancia y susceptancia & \si{\siemens} \\
$S = P + \jj Q$ & Potencia compleja & \si{\volt\ampere} \\
$P$ & Potencia activa (la que realmente se consume o transmite) & \si{\watt} \\
$Q$ & Potencia reactiva (energía que va y vuelve) — no confundir con el factor de calidad & var \\
$f_0$ & Frecuencia de resonancia & \si{\hertz} \\
$Q$ (factor) & Factor de calidad: relación entre energía almacenada y disipada por ciclo & --- \\
$BW$ & Ancho de banda & \si{\hertz} \\
\bottomrule
\end{longtable}

\section*{Líneas de transmisión y ondas}
\begin{longtable}{@{}>{\centering\arraybackslash}p{2.6cm}p{9.6cm}p{2.4cm}@{}}
\toprule \textbf{Símbolo} & \textbf{Significado} & \textbf{Unidad}\\ \midrule\endhead
$c$ & Velocidad de la luz en el vacío, $\num{299792458}$ & \si{\metre\per\second} \\
$\lambda$, $\lambda_0$ & Longitud de onda en la línea y en el vacío & \si{\metre} \\
$v_p$ & Velocidad de fase de la onda en la línea & \si{\metre\per\second} \\
$v_f = v_p/c$ & Factor de velocidad & --- \\
$\varepsilon_r$, $\varepsilon_{\text{eff}}$ & Permitividad relativa del dieléctrico y efectiva (microstrip) & --- \\
$R', L', G', C'$ & Parámetros distribuidos: resistencia, inductancia, conductancia y capacidad \emph{por metro} & \si{\ohm\per\metre}, \si{\henry\per\metre}, \si{\siemens\per\metre}, \si{\farad\per\metre} \\
$Z_0$ & Impedancia característica de la línea (o impedancia de referencia) & \si{\ohm} \\
$\gamma = \alpha + \jj\beta$ & Constante de propagación & \si{\per\metre} \\
$\alpha$ & Constante de atenuación & \si{\neper\per\metre} \\
$\beta = 2\pi/\lambda$ & Constante de fase (número de onda) & \si{\radian\per\metre} \\
$\ell$ & Longitud física de la línea & \si{\metre} \\
$d$ & Distancia medida desde la carga hacia el generador & \si{\metre} \\
$\theta = \beta\ell$ & Longitud eléctrica & \si{\radian} o $^\circ$ \\
$V^+$, $V^-$ & Amplitudes de la onda incidente y de la reflejada & \si{\volt} \\
$Z_L$, $Z_{in}$ & Impedancia de carga y de entrada & \si{\ohm} \\
$\delta_s$ & Profundidad de penetración (efecto pelicular) & \si{\metre} \\
$R_s$ & Resistencia superficial del conductor & \si{\ohm} \\
\bottomrule
\end{longtable}

\section*{Reflexión, adaptación y redes}
\begin{longtable}{@{}>{\centering\arraybackslash}p{2.6cm}p{9.6cm}p{2.4cm}@{}}
\toprule \textbf{Símbolo} & \textbf{Significado} & \textbf{Unidad}\\ \midrule\endhead
$\Gamma$ & Coeficiente de reflexión (complejo) & --- \\
$\Gamma_L$, $\Gamma_{in}$ & Coeficiente de reflexión en la carga y en la entrada & --- \\
$z = Z/Z_0 = r + \jj x$ & Impedancia normalizada (sin unidades) & --- \\
$y = 1/z = g + \jj b$ & Admitancia normalizada & --- \\
$\mathrm{VSWR}$ (ROE) & Relación de onda estacionaria & --- \\
$RL$ & Return Loss, pérdida de retorno (positiva) & \si{\decibel} \\
$ML$ & Mismatch loss, pérdida por desadaptación & \si{\decibel} \\
$Z_s$ & Impedancia de la fuente & \si{\ohm} \\
$\Gamma_p$ & Coeficiente de reflexión de onda de potencia (para fuentes complejas) & --- \\
$Q_L$ & Factor de calidad de la carga, $\abs{X_L}/R_L$ & --- \\
$\Gamma_m$ & Máxima reflexión admitida en la banda de adaptación & --- \\
$G_T$ & Ganancia de transductor & --- o \si{\decibel} \\
$A, B, C, D$ & Parámetros de la matriz de transmisión (cascada) & mixtas \\
\bottomrule
\end{longtable}

\section*{Parámetros S, estabilidad y antenas}
\begin{longtable}{@{}>{\centering\arraybackslash}p{2.6cm}p{9.6cm}p{2.4cm}@{}}
\toprule \textbf{Símbolo} & \textbf{Significado} & \textbf{Unidad}\\ \midrule\endhead
$a_i$, $b_i$ & Ondas de potencia incidente y reflejada en el puerto $i$ & $\sqrt{\mathrm{W}}$ \\
$S_{ij}$ & Parámetro S: respuesta en el puerto $i$ a una excitación en el puerto $j$ & --- \\
$\Delta$ & Determinante de la matriz S de dos puertos & --- \\
$K$, $\mu$ & Factores de estabilidad de Rollett y de Edwards–Sinsky & --- \\
$\tau_g$ & Retardo de grupo & \si{\second} \\
$S_{dd}$, $S_{cc}$, $S_{dc}$ & Parámetros en modo mixto: diferencial, común y conversión & --- \\
$e_{00}$, $e_{11}$, $e_{10}e_{01}$ & Términos de error de la calibración de un puerto & --- \\
$\eta_0$ & Impedancia del espacio libre, $\approx \SI{376.73}{\ohm}$ & \si{\ohm} \\
$k = 2\pi/\lambda$ & Número de onda en el espacio libre & \si{\radian\per\metre} \\
$F(\theta)$ & Diagrama de campo (dependencia angular) & --- \\
$R_r$, $R_L$ & Resistencia de radiación y de pérdidas de la antena & \si{\ohm} \\
$D$, $G$, $e$ & Directividad, ganancia y eficiencia de radiación & ---, \si{\dBi} \\
$\sigma$, $\mu_0$ & Conductividad del metal y permeabilidad del vacío ($4\pi\times10^{-7}$) & \si{\siemens\per\metre}, \si{\henry\per\metre} \\
\bottomrule
\end{longtable}

\section*{Prefijos del Sistema Internacional}
La aplicación acepta estos prefijos al escribir valores: \texttt{10n} = \SI{10}{\nano\henry} si es una bobina, \texttt{4.7p} = \SI{4.7}{\pico\farad}, \texttt{433MHz}, \texttt{2.4G}.
\begin{center}
\begin{tabular}{cccccccccc}\toprule
f & p & n & µ (u) & m & --- & k & M & G & T\\\midrule
$10^{-15}$ & $10^{-12}$ & $10^{-9}$ & $10^{-6}$ & $10^{-3}$ & $1$ & $10^{3}$ & $10^{6}$ & $10^{9}$ & $10^{12}$\\\bottomrule
\end{tabular}
\end{center}
